\section{Problem formulation}
\label{sec:problem-formulation}

% WRITING PROMPTS (comments only; they do not print). About half a page.
% Use notes/method-draft.tex for the notation; keep the same symbols.
%
% Inputs: the lawn polygon L with obstacle holes, the cutting width w, the
% robot start pose, and the robot model (Section robot-energy-model).
% Output: one path from the start pose that covers the free area.
% Decision variables: the decomposition direction theta, the cell order pi,
% and the variant (sweep angle and start corner) a_k of each cell.
% Objective: the battery energy J(theta, pi, a) of the whole path, with
% straights and point turns (Eq. path-energy and Eq. plan-cost in the draft).
% Constraints: the path stays in the free area (lawn eroded by w/2), every
% cell is covered once.
%
% Say explicitly what is NOT optimised (the A* transitions are shortest in
% length, turns are in place, no arc turns) so a reviewer sees the scope.
% State the size of the search space: n! orders times K^n variants, and say
% why the variant choice is exact (dynamic programming, O(n K^2)).
